\section{Why the Batching Ceiling Is Likely Architectural}
\label{app:batching_architectural}

The diffusion batching ceiling ($\leq 1.7\times$) is measured on two
engine-bounded stacks (Appendix~\ref{app:batching}). We argue here that
the ceiling is likely architectural, not an engine bug, by reasoning
about the mechanism rather than measuring it.

\paragraph{The KV-cache reuse argument.}
Autoregressive batching is cheap because of KV-cache reuse: when two
requests share a prompt prefix, the prefill is computed once and the KV
cache is shared. Even without prefix sharing, each request's KV cache
persists across decode steps, so the GPU alternates between requests at
low marginal cost per token.

Block-diffusion invalidates both forms of reuse:

\begin{enumerate}[nosep]
\item \textbf{Cross-request prefix sharing} requires that both requests
are at the same denoising step. vLLM's mixed-phase crash is an engine
bug, but even if it were fixed, requests at different $K$ steps cannot
share a KV cache because the attention mask differs per step.
\item \textbf{Within-request KV persistence} across denoising steps
requires that the keys and values for prompt tokens are identical across
steps. They are --- prompt tokens are never masked --- so in principle
the prompt KV \emph{could} be cached. But the denoising tokens' KV changes
at every step (the mask changes), so the per-step computation cannot
reuse the previous step's activations for the denoising positions.
\end{enumerate}

The net effect: each denoising step is a fresh forward pass over the full
sequence (prompt + denoising positions), with a different attention mask.
The GPU does the same FLOPs whether batching or not, because the
attention pattern is request-specific and step-specific. This is
structurally different from AR decoding, where the attention pattern is
causal and the KV cache grows monotonically.

\paragraph{The cache-off test confirms the direction.}
Disabling prefix caching on the AR arm drops its multiplier from
$42.8\times$ to $26.8\times$ --- still $20\times$ above the DG ceiling.
Cache reuse accounts for part of the AR advantage but not all of it; the
remainder is the bandwidth-bound decode loop, which DG also lacks because
each step is a full forward pass.

\paragraph{What a fixed engine would change.}
If vLLM's mixed-phase crash were fixed (phase-segregated batching), the
DG multiplier might rise above $1.7\times$ --- but only to the extent that
requests at the same $K$ step can share prompt prefill. The per-step
denoising computation remains request-specific. The ceiling would likely
move from $\sim 1\times$ to $\sim 2$--$3\times$, still an order of
magnitude below AR's $26$--$43\times$. The ``true ceiling is unknown''
framing in Appendix~\ref{app:batching} is honest; this appendix adds the
structural argument for why it is unlikely to be high.

\paragraph{Caveat.}
This is a structural argument, not a measurement. A patched vLLM or a
third diffusion serving stack could produce a different number. The
argument says the ceiling is unlikely to reach AR's level, not that it
cannot improve at all.
